% =====================================================================
%  Chương III — PHƯƠNG PHÁP THỰC HIỆN
%  File này được thiết kế để \input vào file luận văn chính.
%  Gói cần khai báo ở preamble của main.tex (đã dùng ở Chương II):
%     \usepackage{amsmath, amssymb}
%     \usepackage{bbm}            % cho \mathbbm{1}
%     \usepackage{booktabs}
%     \usepackage{float}          % cho [H]
%     \usepackage{hyperref}
%     \usepackage{tikz}
%     \usetikzlibrary{arrows.meta, positioning, fit, backgrounds}
%  Nội dung được đối chiếu trực tiếp với mã nguồn trong repo:
%     dataset_utils.py                     (parser .apc, xây dựng lcf_vec, supplement)
%     models/fast_lcf_bert_multitask.py    (kiến trúc APC: LCF + ToMe + đa nhiệm)
%     token_merging/tome_1d.py             (ba chiến lược Token Merging)
%     experiments/run_joint_experiments.py (huấn luyện đa nhiệm, CONFIGS, loss)
%     experiments/eval_joint_triplet.py    (đánh giá bộ ba, Joint/Pipeline)
%     src/ , gas/                          (ATE bằng T5)
% =====================================================================
% Ký hiệu có/không trong bảng cấu hình (dùng amssymb cho \checkmark).
\providecommand{\cmark}{\ensuremath{\checkmark}}
\providecommand{\xmark}{\ensuremath{\times}}

\chapter{III. PHƯƠNG PHÁP THỰC HIỆN}
\label{chap:phuongphap}

Chương này trình bày cách hệ thống được thiết kế và hiện thực. Mục~\ref{sec:overall}
mô tả kiến trúc tổng thể dưới dạng một chuỗi các \emph{stage} xử lý, giải
thích vai trò của từng stage và đi sâu vào chi tiết stage phân loại.
Mục~\ref{sec:methods} trình bày hai phương pháp kỹ thuật cốt lõi được
tích hợp vào kiến trúc -- cơ chế tập trung ngữ cảnh cục bộ (LCF) và kỹ
thuật gộp token (Token Merging). Các mục còn lại trình bày bộ dữ liệu,
hàm mất mát, thiết lập huấn luyện và phương pháp đánh giá. Mọi mô tả đều
được đối chiếu trực tiếp với mã nguồn trong repo.

% =====================================================================
\section{Kiến trúc tổng thể}
\label{sec:overall}

\subsection{Bài toán và luồng xử lý đầu-cuối}

Hệ thống giải quyết bài toán \textbf{trích xuất bộ ba}
(\textit{Joint Triplet Extraction}): từ một câu đánh giá khách sạn $x$,
trích xuất tập

\begin{equation}
\mathcal{T} = \{(a_i, c_i, s_i)\}_{i=1}^{m}
\label{eq:triplet}
\end{equation}

với $a_i$ là \textit{aspect term} (cụm từ khía cạnh), $c_i$ là
\textit{aspect category} (\texttt{AMENITY}, \texttt{BRANDING},
\texttt{EXPERIENCE}, \texttt{FACILITY}, \texttt{LOYALTY},
\texttt{SERVICE}) và $s_i$ là cực tính (\textit{Positive} /
\textit{Negative} / \textit{Neutral}). Một bộ ba được tính đúng khi và
chỉ khi \textbf{cả ba thành phần} khớp nhãn gold.

Hệ thống được tổ chức thành bốn stage nối tiếp
(Hình~\ref{fig:overall}): tiền xử lý, trích xuất khía cạnh (ATE), phân
loại nhóm khía cạnh và cảm xúc (APC), và tổng hợp bộ ba.

\begin{figure}[H]
\centering
\begin{tikzpicture}[
  node distance=0.5cm and 1.0cm,
  stage/.style={rectangle, draw, rounded corners, align=center, fill=black!4,
                minimum height=1.0cm, minimum width=8.0cm, font=\small},
  io/.style={rectangle, draw, dashed, align=center, rounded corners,
             minimum height=0.8cm, minimum width=6.0cm, font=\small\itshape},
  arr/.style={-{Stealth[length=2.2mm]}, thick}
]
  \node[io] (in) {Câu đánh giá khách sạn (input)};
  \node[stage, below=of in] (s1)
       {\textbf{Stage 1 -- Tiền xử lý \& biểu diễn đầu vào}\\
        Token hóa SPC \texttt{[CLS] câu [SEP] aspect [SEP]} + dựng vector \texttt{lcf\_vec}};
  \node[stage, below=of s1] (s2)
       {\textbf{Stage 2 -- ATE (T5)}\\
        Sinh có điều kiện danh sách aspect term + chuẩn hóa Levenshtein};
  \node[stage, below=of s2] (s3)
       {\textbf{Stage 3 -- APC (\texttt{FastLcfBertMultiTask})}\\
        Backbone $\to$ Token Merging $\to$ LCF $\to$ phân loại đa nhiệm};
  \node[stage, below=of s3] (s4)
       {\textbf{Stage 4 -- Tổng hợp bộ ba}\\
        Ghép $(a_i, c_i, s_i)$};
  \node[io, below=of s4] (out) {Tập bộ ba $\mathcal{T} = \{(a_i, c_i, s_i)\}$ (output)};

  \draw[arr] (in) -- (s1);
  \draw[arr] (s1) -- (s2);
  \draw[arr] (s2) -- (s3) node[midway, right, font=\scriptsize] {aspect terms};
  \draw[arr] (s3) -- (s4) node[midway, right, font=\scriptsize] {(category, sentiment)};
  \draw[arr] (s4) -- (out);
\end{tikzpicture}
\caption{Kiến trúc tổng thể của hệ thống trích xuất bộ ba, tổ chức thành
bốn stage nối tiếp. Đường \texttt{aspect terms} từ Stage 2 sang Stage 3
là luồng của kiến trúc Pipeline; ở kiến trúc Joint, aspect term được lấy
trực tiếp từ nhãn gold (mục~\ref{subsec:joint_pipeline}).}
\label{fig:overall}
\end{figure}

\subsection{Giải thích từng stage}

\paragraph{Stage 1 — Tiền xử lý và biểu diễn đầu vào.}
Mỗi mẫu được token hóa theo định dạng \textbf{Sentence-Pair
Classification (SPC)}, ghép câu (segment A) với cụm khía cạnh
(segment B): \texttt{[CLS] <câu> [SEP] <aspect term> [SEP]}, độ dài tối
đa $L = 128$. Đồng thời, hệ thống dựng \textbf{vector chỉ thị khía cạnh}
\texttt{lcf\_vec} $\in \{0,1\}^{L}$: dựa trên \texttt{offset\_mapping}
của tokenizer, mỗi token thuộc segment A có khoảng ký tự \emph{giao} với
khoảng ký tự của cụm khía cạnh được gán 1, còn lại bằng 0. Nếu cụm khía
cạnh bị cắt khỏi segment A, hệ thống quay lui dùng
\texttt{token\_type\_ids} (bản sao khía cạnh ở segment B) làm chỉ thị.
Vector này là tín hiệu đầu vào cho phương pháp LCF (mục~\ref{subsec:lcf}).

\paragraph{Stage 2 — Trích xuất khía cạnh (ATE) bằng T5.}
Bài toán ATE được hình thức hóa thành \textbf{sinh có điều kiện}: mô hình
\texttt{T5AspectExtractor} (backbone \texttt{t5-base}) nhận câu và sinh
trực tiếp danh sách cụm từ khía cạnh theo định dạng GAS, ví dụ
\texttt{"(phòng); (nhân viên); (thang máy)"}; câu không có khía cạnh sinh
ra \texttt{"none"}. Quá trình sinh dùng beam search (beam $=4$,
\texttt{max\_len} $=64$). Vì chuỗi sinh có thể lệch khỏi văn bản gốc, đầu
ra được \textbf{chuẩn hóa bằng đối sánh n-gram theo khoảng cách
Levenshtein}, ánh xạ mỗi cụm từ về chuỗi con gần nhất thực sự xuất hiện
trong câu.

\paragraph{Stage 3 — Phân loại (APC) đa nhiệm.}
Với mỗi cặp (câu, aspect term), mô hình \texttt{FastLcfBertMultiTask}
dự đoán đồng thời nhóm khía cạnh và cảm xúc. Đây là stage trung tâm, nơi
hai phương pháp LCF và Token Merging được tích hợp; chi tiết kiến trúc
được trình bày ở mục~\ref{subsec:apc_detail}.

\paragraph{Stage 4 — Tổng hợp bộ ba.}
Aspect term (từ Stage 2 hoặc nhãn gold) được ghép với nhóm khía cạnh và
cảm xúc (từ Stage 3) thành các bộ ba $(a_i, c_i, s_i)$ -- đầu ra cuối
cùng của hệ thống.

\subsection{Hai kiến trúc tổ chức: Joint và Pipeline}
\label{subsec:joint_pipeline}

Khóa luận khảo sát hai cách kết nối Stage 2 và Stage 3
(Bảng~\ref{tab:joint_pipeline}):

\begin{table}[H]
\centering
\caption{Hai kiến trúc tổ chức bài toán đầu-cuối}
\label{tab:joint_pipeline}
\begin{tabular}{@{}p{2.0cm}p{6.5cm}p{4.5cm}@{}}
\toprule
\textbf{Kiến trúc} & \textbf{Mô tả} & \textbf{Đặc điểm} \\
\midrule
\textbf{Joint} & APC được đánh giá trên aspect term \emph{gold}; ATE và
APC không phụ thuộc đầu ra của nhau & Cô lập được năng lực bộ phân loại,
không chịu lan truyền lỗi \\
\textbf{Pipeline} & T5 thực hiện ATE trước, aspect term \emph{dự đoán}
được đưa sang APC & Linh hoạt, phản ánh hiệu suất thực tế; chịu lan
truyền lỗi từ ATE \\
\bottomrule
\end{tabular}
\end{table}

\subsection{Chi tiết kiến trúc Stage APC}
\label{subsec:apc_detail}

Hình~\ref{fig:apc_detail} mô tả chi tiết stage phân loại. Đầu ra
\texttt{last\_hidden\_state} $H \in \mathbb{R}^{L \times d_h}$ của
backbone đi qua khối \textbf{Token Merging} (tùy chọn), sau đó tách thành
\textbf{hai luồng}: luồng cục bộ (nhân với trọng số LCF) và luồng toàn
cục (giữ nguyên). Hai luồng được hợp nhất, đưa qua self-attention và
pooler, rồi vào hai đầu phân loại.

\begin{figure}[H]
\centering
\begin{tikzpicture}[
  node distance=0.5cm and 0.9cm,
  block/.style={rectangle, draw, rounded corners, align=center,
                 minimum height=0.85cm, minimum width=3.4cm, font=\small},
  optblock/.style={block, dashed, fill=black!4},
  head/.style={block, minimum width=3.0cm},
  arr/.style={-{Stealth[length=2.2mm]}, thick}
]
  \node[block] (input) {Input: \texttt{input\_ids},\\\texttt{attention\_mask}, \texttt{lcf\_vec}};
  \node[block, below=of input] (bert) {Backbone (BERT)\\$H \in \mathbb{R}^{L\times d_h}$};
  \node[optblock, below=of bert] (tome) {\textbf{Token Merging} (tùy chọn)\\$\to L' \to$ resize $\to L$};
  \node[block, below=of tome] (lcf) {\textbf{LCF}: $H_{local}=H\odot c$\\(CDW mặc định / CDM)};
  \node[block, below=of lcf] (sa1) {Self-Attention (\texttt{bert\_SA})};
  \node[block, below=of sa1] (fusion) {Fusion: Concat$(H_{local}, H)$\\$\to$ Linear$(2d\!\to\!d)$};
  \node[block, below=of fusion] (sa2) {Self-Attention (\texttt{bert\_SA\_})};
  \node[block, below=of sa2] (pool) {Pooler (token \texttt{[CLS]})};
  \node[head, below left=0.7cm and -1.4cm of pool] (h1) {Sentiment head\\Linear$(d\!\to\!3)$};
  \node[head, below right=0.7cm and -1.4cm of pool] (h2) {Category head\\Linear$(d\!\to\!6)$};

  \draw[arr] (input) -- (bert);
  \draw[arr] (bert) -- (tome);
  \draw[arr] (tome) -- (lcf);
  \draw[arr] (lcf) -- (sa1);
  \draw[arr] (sa1) -- (fusion);
  \draw[arr] (fusion) -- (sa2);
  \draw[arr] (sa2) -- (pool);
  \draw[arr] (pool) -- (h1);
  \draw[arr] (pool) -- (h2);
  % nhánh toàn cục đi vòng tới fusion
  \draw[arr] (tome.east) -- ++(2.9,0) |- (fusion.east);
\end{tikzpicture}
\caption{Chi tiết kiến trúc Stage APC (\texttt{FastLcfBertMultiTask}).
Khối Token Merging (nét đứt) là tùy chọn; biểu diễn toàn cục đi vòng tới
bước Fusion. Hai đầu phân loại chia sẻ chung phần thân.}
\label{fig:apc_detail}
\end{figure}

% =====================================================================
\section{Các phương pháp}
\label{sec:methods}

Hai phương pháp kỹ thuật được tích hợp vào Stage APC nhằm hai mục tiêu
bổ trợ nhau: \textbf{LCF} (mục~\ref{subsec:lcf}) tăng độ chính xác phân
loại bằng cách tập trung vào ngữ cảnh quanh khía cạnh; \textbf{Token
Merging} (mục~\ref{subsec:tome}) giảm chi phí tính toán bằng cách rút
ngắn chuỗi token. Nền tảng lý thuyết của hai phương pháp đã trình bày ở
Chương~\ref{chap:coso_lythuyet}; phần này tập trung vào cách \emph{áp
dụng và hiện thực} chúng trong hệ thống.

\subsection{Phương pháp 1: Tập trung ngữ cảnh cục bộ (LCF)}
\label{subsec:lcf}

\paragraph{Ý tưởng.} Với mỗi khía cạnh, các từ ở gần thường mang thông
tin cảm xúc liên quan hơn các từ ở xa. LCF \cite{zeng2019lcf,
yang2021lcfatepc} khai thác điều này bằng cách tạo thêm một luồng biểu
diễn ``cục bộ'' trong đó token xa khía cạnh bị giảm trọng số, rồi kết
hợp với luồng toàn cục.

\paragraph{Khoảng cách ngữ nghĩa tương đối (SRD).} Từ vector
\texttt{lcf\_vec} (dựng ở Stage 1), hệ thống xác định tâm và nửa độ rộng
của cụm khía cạnh, rồi tính SRD của token tại vị trí $i$:

\begin{equation}
SRD_i = \max\!\Big(0,\; \big|\,i - \text{center}\,\big| - \text{half}\Big),
\qquad
\text{half} = \left\lfloor \frac{a_{end} - a_{start} + 1}{2} \right\rfloor
\label{eq:srd_m}
\end{equation}

với ngưỡng $\alpha = 5$ (\texttt{SRD\_THRESHOLD}): token có $SRD_i \le \alpha$
thuộc ngữ cảnh cục bộ.

\paragraph{Hai biến thể tạo trọng số.} Hệ thống cài đặt hai biến thể,
chọn qua cờ \texttt{use\_cdm}:

\begin{itemize}
  \item \textbf{CDW} (Context Dynamic Weight -- mặc định): trọng số giảm
        dần liên tục, $c_i = 1$ nếu $SRD_i \le \alpha$, ngược lại
        $c_i = \big(L - (SRD_i - \alpha)\big)/L$ (kẹp về $[0,1]$);
  \item \textbf{CDM} (Context Dynamic Mask): mặt nạ nhị phân,
        $m_i = 1$ nếu $SRD_i \le \alpha$, ngược lại $m_i = 0$.
\end{itemize}

\paragraph{Tích hợp hai luồng.} Luồng cục bộ $H_{local} = H \odot c$
(hoặc $H \odot m$) qua self-attention \texttt{bert\_SA}; luồng toàn cục
$H$ giữ nguyên. Hai luồng được hợp nhất:

\begin{equation}
H_{fused} = \texttt{Linear}_{2d \rightarrow d}\big(\,[\,H_{local}\,;\,H\,]\,\big)
\label{eq:fusion_m}
\end{equation}

trước khi qua self-attention thứ hai và pooler (xem
Hình~\ref{fig:apc_detail}).

\subsection{Phương pháp 2: Gộp token (Token Merging -- ToMe)}
\label{subsec:tome}

\paragraph{Ý tưởng.} Chi phí self-attention tỉ lệ $O(n^2)$ theo độ dài
chuỗi $n$; nhiều token (từ nối, dấu câu, subword lặp) mang ít thông tin
phân biệt cảm xúc. ToMe \cite{bolya2023tome} gộp các token có biểu diễn
tương tự nhau (đo bằng cosine) thành một token đại diện bằng phép trung
bình, rút ngắn chuỗi mà không thêm tham số. Trong hệ thống, ToMe được áp
dụng \textbf{sau backbone, trước LCF} (module
\texttt{ToMeSequenceMerger}).

\paragraph{Ba chiến lược gộp.} Lựa chọn qua tham số
\texttt{merge\_strategy}:

\begin{table}[H]
\centering
\caption{Ba chiến lược gộp token được khảo sát}
\label{tab:tome_strategies}
\begin{tabular}{@{}p{3.4cm}p{10.4cm}@{}}
\toprule
\textbf{Chiến lược} & \textbf{Cơ chế} \\
\midrule
\texttt{bipartite} & Chia token nội bộ thành hai nhóm xen kẽ $A$, $B$;
mỗi token $A$ ghép với token $B$ tương đồng cosine nhất (phương pháp gốc,
CVPR 2023). \\
\texttt{sequential\_local} & Quét trái$\to$phải, mỗi token gộp về phía
hàng xóm (trái/phải) giống nhất; giữ nguyên ranh giới vào vùng khía cạnh. \\
\texttt{attention\_weighted} & Ưu tiên gộp token có trọng số attention
thấp nhất; bảo vệ token khía cạnh và token attention cao (top 25\%). \\
\bottomrule
\end{tabular}
\end{table}

\paragraph{Bảo vệ token quan trọng.} Cả ba chiến lược đều bảo vệ tuyệt
đối \texttt{[CLS]}, \texttt{[SEP]} và token thuộc vùng khía cạnh
(\texttt{protect\_aspect=True}). Khi xảy ra gộp, chỉ thị khía cạnh được
lan truyền qua phép lấy cực đại: $\text{lcf}[\text{src}] =
\max(\text{lcf}[\text{src}], \text{lcf}[\text{dst}])$ -- nhờ đó LCF ở
bước sau vẫn hoạt động đúng trên chuỗi đã rút gọn.

\paragraph{Khôi phục độ dài (resize).} Sau khi gộp, chuỗi có độ dài
$L' \le L$. Trong khóa luận, kết quả gộp \textbf{luôn được nội suy tuyến
tính trở lại độ dài gốc $L$} (\texttt{resize=True}, dùng
\texttt{F.interpolate} với \texttt{mode="linear"}) cho cả hidden states và
\texttt{lcf\_vec}. Nhờ đó các lớp phía sau (LCF, self-attention, pooler)
giữ nguyên kiến trúc và độ dài đầu vào không đổi qua mọi cấu hình.

\subsubsection{Vì sao chỉ dùng resize mà không dùng compact}
\label{subsubsec:why_resize}

Module \texttt{ToMeSequenceMerger} có hỗ trợ chế độ \emph{compact}
(\texttt{resize=False}) -- giữ nguyên chuỗi rút gọn độ dài $L'$ và đệm
theo $L'_{\max}$ của batch. Tuy nhiên, \textbf{toàn bộ thực nghiệm của
khóa luận chỉ dùng \texttt{resize=True}}, vì ba lý do sau:

\begin{enumerate}
  \item \textbf{Cô lập đúng biến cần đo.} Câu hỏi nghiên cứu đối với ToMe
        là \emph{liệu việc gộp token dư thừa có làm suy giảm chất lượng
        phân loại không, và chiến lược gộp nào bảo toàn thông tin tốt
        nhất} -- một câu hỏi về \emph{chất lượng biểu diễn}. Chế độ resize
        giữ độ dài chuỗi cố định ở $L = 128$ qua mọi cấu hình, nên khác
        biệt duy nhất giữa baseline và một cấu hình có gộp là \emph{nội
        dung} biểu diễn (do trung bình các token tương đồng). Mọi thay đổi
        về độ chính xác vì thế quy hoàn toàn về thao tác gộp, không bị
        nhiễu bởi độ dài chuỗi hay cách đệm khác nhau.
  \item \textbf{Compact gây nhiễu so sánh.} Ở chế độ compact, mỗi cấu hình
        (và mỗi câu) sinh ra độ dài $L'$ khác nhau, buộc phải đệm theo
        $L'_{\max}$ của batch. Điều này trộn lẫn ảnh hưởng của \emph{việc
        gộp} với ảnh hưởng của \emph{độ dài chuỗi thay đổi}, khiến không
        thể quy kết sạch chênh lệch độ chính xác cho riêng chiến lược gộp.
  \item \textbf{Lợi ích tốc độ của compact ở vị trí này là không đáng kể.}
        ToMe được chèn \emph{sau} backbone (mục~\ref{subsec:tome}), nên 12
        lớp Transformer của backbone -- phần tốn kém nhất -- vẫn luôn chạy
        ở độ dài đầy đủ $L = 128$ bất kể chế độ nào. Compact chỉ rút ngắn
        đầu vào cho hai lớp self-attention nhẹ phía sau, vốn chiếm tỉ lệ
        nhỏ trong tổng chi phí. Do đó compact gần như không mang lại tăng
        tốc thực sự, trong khi lại làm phức tạp việc so sánh.
\end{enumerate}

\noindent Vì những lý do trên, khóa luận giới hạn phạm vi thực nghiệm ở
\textbf{câu hỏi chất lượng biểu diễn} (dùng resize). Câu hỏi \emph{tăng
tốc suy luận} thực sự -- vốn đòi hỏi gộp token \emph{trước} backbone
(pre-BERT ToMe, để chuỗi ngắn đi vào ngay 12 lớp Transformer) hoặc dùng
compact ở vị trí thích hợp -- được để lại cho hướng phát triển
(Chương~5).

\subsection{Huấn luyện đa nhiệm}
\label{subsec:multitask_loss}

Hai đầu phân loại được huấn luyện đồng thời bằng tổng hai thành phần
\texttt{CrossEntropyLoss}:

\begin{equation}
\mathcal{L} = w_{s}\,\mathcal{L}_{\text{sentiment}}
            + w_{c}\,\mathbbm{1}[\text{main\_mask}]\cdot\mathcal{L}_{\text{aspect\_cat}}
\label{eq:loss_m}
\end{equation}

trong đó $\mathbbm{1}[\text{main\_mask}]$ là hàm chỉ thị chỉ bật trên mẫu
chính (\texttt{is\_supplement=False}) -- đầu phân loại nhóm khía cạnh
\textbf{không} học trên mẫu bổ trợ (vốn không có nhãn category). Cấu hình
chính dùng $w_s = w_c = 1.0$. Cả hai thành phần dùng \textbf{trọng số lớp
theo tần suất nghịch đảo} $w_k = N/(C \cdot n_k)$ ($N$: tổng số mẫu, $C$:
số lớp, $n_k$: số mẫu lớp $k$) để giảm ảnh hưởng của mất cân bằng dữ liệu.
Giá trị siêu tham số và quy trình huấn luyện cụ thể được trình bày ở
Chương~4.

% =====================================================================
\section{Ví dụ minh họa chi tiết các phương pháp}
\label{sec:examples}

Để làm rõ cơ chế của hai phương pháp, mục này trình bày một \textbf{ví dụ
chạy xuyên suốt} với các vector ``đồ chơi'' (toy vectors) 4 chiều thay cho
hidden states $768$ chiều thực tế. Các con số được chọn để minh họa, song
mọi công thức và quy tắc đều đúng với mã nguồn.

\subsection{Thiết lập chung}
\label{subsec:ex_setup}

Xét câu \textit{``I found the food delicious but service was terrible''}
với khía cạnh \textbf{food} (gold: \textit{positive}, \texttt{FOOD}). Sau
token hóa SPC, chuỗi gồm 13 token hợp lệ
(\texttt{[CLS]} + câu + \texttt{[SEP]} + \texttt{food} + \texttt{[SEP]}),
phần còn lại là padding. Bảng~\ref{tab:ex_tokens} liệt kê vị trí, token,
vector ẩn $v_i$ và giá trị \texttt{lcf\_vec} (chỉ token khía cạnh ở
segment A được gán 1).

\begin{table}[H]
\centering
\caption{Thiết lập ví dụ: 13 token hợp lệ với vector ẩn 4 chiều}
\label{tab:ex_tokens}
\small
\begin{tabular}{@{}clcc@{}}
\toprule
\textbf{Vị trí} & \textbf{Token} & \textbf{Vector $v_i$} & \texttt{lcf\_vec} \\
\midrule
0  & \texttt{[CLS]}      & $[1,0,1,0]$ & 0 \\
1  & i                   & $[1,1,0,1]$ & 0 \\
2  & found               & $[2,1,2,1]$ & 0 \\
3  & the                 & $[1,0,1,1]$ & 0 \\
4  & \textbf{food} (aspect) & $[5,4,3,2]$ & \textbf{1} \\
5  & delicious           & $[4,3,2,2]$ & 0 \\
6  & but                 & $[1,0,2,0]$ & 0 \\
7  & service             & $[3,2,1,2]$ & 0 \\
8  & was                 & $[1,1,2,0]$ & 0 \\
9  & terrible            & $[3,1,0,3]$ & 0 \\
10 & \texttt{[SEP]}$_{\text{mid}}$ & $[1,0,2,1]$ & 0 \\
11 & food$_B$ (bản sao)  & $[5,4,3,2]$ & 0 \\
12 & \texttt{[SEP]}$_{\text{end}}$ & $[0,1,0,1]$ & 0 \\
\bottomrule
\end{tabular}
\end{table}

\subsection{Ví dụ Phương pháp 1a -- LCF-CDW}
\label{subsec:ex_cdw}

Khía cạnh nằm tại $a_{start}=a_{end}=4$, nên
$\text{center}=4$, $\text{half}=\lfloor 1/2 \rfloor = 0$, và
$SRD_i = |i-4|$. Với $\alpha = 5$ và $L = 128$, trọng số CDW là 1 cho mọi
token có $|i-4| \le 5$ (vị trí 0--9) và giảm nhẹ cho các vị trí xa hơn
(Bảng~\ref{tab:ex_cdw}).

\begin{table}[H]
\centering
\caption{Trọng số CDW cho các vị trí có $SRD > \alpha$}
\label{tab:ex_cdw}
\small
\begin{tabular}{@{}cccl@{}}
\toprule
\textbf{Vị trí} & $SRD_i$ & $c_i = \frac{L-(SRD_i-\alpha)}{L}$ & \textbf{Kết quả $v_i \cdot c_i$} \\
\midrule
0--9 & $\le 5$ & $1.000$ & giữ nguyên \\
10 & 6 & $127/128 = 0.992$ & $[0.99,\,0,\,1.98,\,0.99]$ \\
11 & 7 & $126/128 = 0.984$ & $[4.92,\,3.94,\,2.95,\,1.97]$ \\
12 & 8 & $125/128 = 0.977$ & $[0,\,0.98,\,0,\,0.98]$ \\
\bottomrule
\end{tabular}
\end{table}

\noindent\textbf{Nhận xét:} vì $L=128$ lớn, CDW chỉ giảm trọng số rất nhẹ
(còn $\approx 0.98$) ngay cả với token xa nhất -- token xa vẫn đóng góp,
chỉ ít hơn. Với câu dài hơn (ví dụ 23 token, khía cạnh ở đầu câu), các
token cuối có thể giảm xuống $\approx 0.89$: CDW \emph{giảm mượt} theo
khoảng cách.

\subsection{Ví dụ Phương pháp 1b -- LCF-CDM}
\label{subsec:ex_cdm}

CDM dùng cùng SRD nhưng cắt cứng: $m_i = 1$ nếu $SRD_i \le 5$, ngược lại
$m_i = 0$ (không phụ thuộc $L$). Trên ví dụ, các vị trí 10--12 bị
\textbf{xóa hoàn toàn} (Bảng~\ref{tab:ex_cdm_vs_cdw}).

\begin{table}[H]
\centering
\caption{So sánh CDW và CDM tại các vị trí 10--12}
\label{tab:ex_cdm_vs_cdw}
\small
\begin{tabular}{@{}clll@{}}
\toprule
\textbf{Vị trí} & \textbf{Token} & \textbf{CDW} & \textbf{CDM} \\
\midrule
10 & \texttt{[SEP]}$_{\text{mid}}$ & $[0.99,0,1.98,0.99]$ & $[0,0,0,0]$ \\
11 & food$_B$                       & $[4.92,3.94,2.95,1.97]$ & $[0,0,0,0]$ \\
12 & \texttt{[SEP]}$_{\text{end}}$ & $[0,0.98,0,0.98]$ & $[0,0,0,0]$ \\
\bottomrule
\end{tabular}
\end{table}

\noindent\textbf{Nhận xét:} với câu ngắn, chênh lệch giữa CDM và CDW nhỏ
(chỉ 3 token cuối). Với câu dài, CDM có thể cắt hơn 15 token -- luồng cục
bộ chỉ còn đúng vùng câu quanh khía cạnh, bỏ hẳn \texttt{[SEP]} và bản sao
segment B. Đây là lý do CDM thường tốt hơn cho nhãn thiểu số (tập trung
``cứng''), còn CDW giữ nhiều thông tin toàn cục hơn.

\subsection{Ví dụ Phương pháp 2a -- ToMe Bipartite}
\label{subsec:ex_bip}

Token được bảo vệ: \texttt{[CLS]} (vị trí 0), \texttt{[SEP]}$_{\text{end}}$
(vị trí 12) và token khía cạnh \textbf{food} (vị trí 4). Pool nội bộ còn
lại gồm 10 vị trí $\{1,2,3,5,6,7,8,9,10,11\}$. Chia xen kẽ \emph{theo thứ
hạng trong pool}\footnote{Theo \texttt{torch.arange(0,nm,2)} và
\texttt{arange(1,nm,2)} trong \texttt{\_bipartite\_pairs} -- chia theo
thứ hạng, không theo chẵn/lẻ của vị trí token.}:

\begin{equation*}
A = \{1, 3, 6, 8, 10\}, \qquad B = \{2, 5, 7, 9, 11\}.
\end{equation*}

Mỗi token nhóm $A$ chọn token nhóm $B$ có cosine lớn nhất. Ví dụ với vị
trí 1 (\textit{i}) và vị trí 7 (\textit{service}):

\begin{equation*}
\cos(v_1, v_7) = \frac{[1,1,0,1]\cdot[3,2,1,2]}{\sqrt{3}\,\sqrt{18}}
= \frac{7}{7.35} \approx 0.95.
\end{equation*}

Bảng~\ref{tab:ex_bip} trình bày kết quả chọn cặp và khử trùng (mỗi token
$B$ chỉ được ghép một lần, duyệt $A$ theo thứ tự).

\begin{table}[H]
\centering
\caption{Bipartite -- chọn cặp và khử trùng (bước 1)}
\label{tab:ex_bip}
\small
\begin{tabular}{@{}llcl@{}}
\toprule
\textbf{Token $A$} & \textbf{$B$ tương đồng nhất} & \textbf{cos} & \textbf{Kết quả} \\
\midrule
1 (i)        & 7 (service) & 0.95 & \cmark{} ghép $1\!\leftarrow\!7$ \\
3 (the)      & 2 (found)   & 0.91 & \cmark{} ghép $3\!\leftarrow\!2$ \\
6 (but)      & 2 (found)   & 0.85 & \xmark{} ($B{=}2$ đã dùng) \\
8 (was)      & 2 (found)   & 0.90 & \xmark{} ($B{=}2$ đã dùng) \\
10 (\texttt{[SEP]}) & 2 (found) & 0.90 & \xmark{} ($B{=}2$ đã dùng) \\
\bottomrule
\end{tabular}
\end{table}

Hai cặp được gộp bằng trung bình: vị trí 1 hấp thụ vị trí 7
($x_1 \leftarrow \tfrac12(v_1+v_7)$), vị trí 3 hấp thụ vị trí 2; hai token
2 và 7 bị loại $\Rightarrow$ chuỗi còn \textbf{11 token}. Vì
\texttt{TOME\_MERGE\_STEPS = 2}, quá trình lặp lại trên 11 token và rút
xuống \textbf{$\approx 9$ token}. Chỉ thị khía cạnh được lan truyền qua
phép $\max$ (ở đây cả hai cặp đều có \texttt{lcf}$=0$ nên không đổi).

\subsection{Ví dụ Phương pháp 2b -- ToMe Sequential Local}
\label{subsec:ex_seq}

Chiến lược này quét trái$\to$phải, mỗi token gập về phía hàng xóm giống
nhất. Trước hết, ranh giới vào vùng khía cạnh được phát hiện: tại vị trí
3 có $\texttt{lcf}_3{=}0$ và $\texttt{lcf}_4{=}1$ nên vị trí 3
(\textit{the}) được đánh dấu \texttt{mid\_sep} và \emph{không} bị gộp.

Khi quét tới vị trí 5 (\textit{delicious}):

\begin{align*}
\cos(v_5, v_4) &= \frac{[4,3,2,2]\cdot[5,4,3,2]}{\sqrt{33}\,\sqrt{54}} = \frac{42}{42.2} \approx 0.995 \quad (\text{trái}),\\
\cos(v_5, v_6) &= \frac{[4,3,2,2]\cdot[1,0,2,0]}{\sqrt{33}\,\sqrt{5}} = \frac{8}{12.8} \approx 0.62 \quad (\text{phải}).
\end{align*}

Vì sim trái lớn hơn, \textit{delicious} gập vào \textbf{food} (vị trí 4):
$x_4 \leftarrow \tfrac12(v_4+v_5)$ và quan trọng là
$\texttt{lcf}_4 = \max(1,0) = 1$ -- \textbf{tín hiệu khía cạnh được giữ
nguyên}. Token khía cạnh không bị xóa nhưng \emph{có thể nhận} các token
gộp vào. Do hiệu ứng \textbf{xếp tầng (cascade)} trong một lượt quét, các
token kế tiếp (\textit{but}, \textit{service}, \textit{was},
\textit{terrible}, \dots) lần lượt gập vào vị trí 4, khiến chuỗi rút mạnh
xuống $\approx 5$ token chỉ sau \textbf{một lượt}. Sequential do đó gộp
mạnh hơn bipartite (không bị giới hạn bởi số bước).

\subsection{Ví dụ Phương pháp 2c -- ToMe Attention-Weighted}
\label{subsec:ex_attn}

Trong huấn luyện, vòng lặp không truyền trọng số attention nên
\texttt{attn\_seg} đồng nhất ($=1$). Khi đó thuật toán chọn token
\emph{chưa bị loại, không được bảo vệ, không phải khía cạnh} theo thứ tự,
rồi gộp nó vào hàng xóm cosine cao nhất. Ví dụ ở vòng 1, vị trí 1
(\textit{i}) được chọn; cosine cao nhất là với vị trí 7 (\textit{service},
$\approx 0.95$) nên \textit{i} bị xóa và
$x_7 \leftarrow \tfrac12(v_7+v_1)$ (xóa $i$, cập nhật $j$ bằng trung bình).
Giới hạn $\lfloor(13-2)/2\rfloor = 5$ lần gộp, chuỗi rút xuống
\textbf{8 token}; token khía cạnh \textbf{food} luôn được bảo vệ.

\subsection{Ví dụ tổ hợp: LCF-CDM + Bipartite + resize}
\label{subsec:ex_combined}

Cuối cùng, ghép các bước thành luồng đầy đủ của một cấu hình (ví dụ
\texttt{use\_lcf=True}, \texttt{use\_cdm=True}, \texttt{use\_tome=True},
\texttt{tome\_resize=True}, \texttt{merge\_strategy="bipartite"}):

\begin{enumerate}
  \item \textbf{Backbone} sinh $H \in \mathbb{R}^{128 \times d_h}$ (13
        hàng hợp lệ, còn lại padding).
  \item \textbf{ToMe bipartite} rút 13 $\to 9$ token, rồi
        \textbf{resize} nội suy tuyến tính trở lại 128 hàng. Vector
        \texttt{lcf\_vec} cũng được nội suy nên không còn nhị phân: vùng
        khía cạnh ``loang'' thành một dải quanh vị trí tương ứng.
  \item \textbf{CDM} tính lại từ \texttt{lcf\_vec} đã resize: tìm vùng
        $\texttt{lcf} > 0.5$, suy ra center/half và mặt nạ; chỉ một cửa sổ
        cục bộ quanh khía cạnh được giữ ($m_i = 1$), phần còn lại bằng 0.
  \item \textbf{Luồng cục bộ} $H_{local} = H \odot m$; \textbf{luồng toàn
        cục} giữ nguyên $H$.
  \item \textbf{Fusion} $\to$ self-attention $\to$ pooler $\to$ hai đầu
        phân loại, cho \texttt{sentiment\_logits} (3 lớp) và
        \texttt{aspect\_cat\_logits} (6 lớp); $\arg\max$ cho
        \textit{positive} và \texttt{FOOD}.
\end{enumerate}

\noindent Hai điểm cần nhấn mạnh:

\begin{itemize}
  \item \textbf{ToMe không tạo đặc trưng mới} -- nó chỉ \emph{giảm số
        token} bằng cách trung bình các token tương đồng; đầu ra của ToMe
        được dùng ngay cho LCF, không qua lớp biến đổi nào khác.
  \item \textbf{Token \texttt{[CLS]} đưa vào pooler là từ biểu diễn đã
        hợp nhất} ($H_{fused}$ sau Fusion + self-attention), \emph{không
        phải} \texttt{[CLS]} thô của backbone -- tức quyết định phân loại
        dựa trên cả luồng cục bộ lẫn toàn cục.
\end{itemize}

% =====================================================================
\section*{Tóm tắt chương}
\addcontentsline{toc}{section}{Tóm tắt chương}

Chương này đã trình bày phương pháp thực hiện theo hai trục chính: (1)
\textbf{kiến trúc tổng thể} gồm bốn stage nối tiếp -- tiền xử lý, ATE
(T5), APC (\texttt{FastLcfBertMultiTask}) và tổng hợp bộ ba -- cùng chi
tiết kiến trúc stage phân loại và hai cách tổ chức Joint/Pipeline; (2)
\textbf{hai phương pháp kỹ thuật} tích hợp vào stage APC: LCF (với SRD và
hai biến thể CDM/CDW) nhằm tập trung ngữ cảnh cục bộ, và Token Merging
(ba chiến lược) nhằm giảm dư thừa token. Chương cũng làm rõ \textbf{lý do
toàn bộ thực nghiệm chỉ dùng chế độ \texttt{resize}} (cô lập câu hỏi chất
lượng biểu diễn) và hàm mất mát đa nhiệm có trọng số lớp. Bộ dữ liệu, thiết
lập huấn luyện, không gian cấu hình và kết quả thực nghiệm được trình bày
ở Chương~4.
